\documentclass[a4paper]{article}

%Пакеты для математических символов:
\usepackage{amsmath} % американское математическое сообщество.
\usepackage{amssymb} % миллион разных значков и готический, ажурный шрифты.
\usepackage{amscd} % диаграммы, графики.
\usepackage{amsthm} % окружения теорем, определений и тд.
\usepackage{physics} % основные физические символы
%\usepackage{latexsym} % треугольники и пьяная стрелка.

%пакеты для шрифтов:
%\usepackage{euscript} % прописной шрифт с завитушками.
\usepackage{MnSymbol} % Значеки доказательства
\usepackage{verbatim} % улучшенный шрифт "пишущей машинки".
%\usepackage{array} % более удобные таблицы.
%\usepackage{multirow} % мультистолбцы в таблицах.
%\usepackage{longtable} % таблицы на несколько страниц.



%Пакеты для оформления:
\RequirePackage[center, medium]{titlesec}% Стиль секций и заголовков
\usepackage{collectbox} % Добавляет коробочки, можно складывать туда текст)
%\usepackage[x11names]{xcolor} % 317 новых цветов для текста.
%\usepackage{multicol} % набор текста в несколько колонн.
\usepackage{graphicx} % расширенные возможности вставки стандартных картинок.
\usepackage{subcaption} % возможность вставлять картинки в строчку
%\usepackage{caption} % возможность подавить нумерацию у caption.
\usepackage{wrapfig} % вставка картинок и таблиц, обтекаемых текстом.
\usepackage{cancel} % значки для сокращения дробей, упрощения, стремления.
\usepackage{misccorr} % в заголовках появляется точка, но при ссылке на них ее нет.
%\usepackage{indentfirst} % отступ у первой строки раздела
%\usepackage{showkeys} % показывает label формул над их номером.
%\usepackage{fancyhdr} % удобное создание верхних и нижних колонтитулов.
%\usepackage{titlesec} % еще одно создание верхних и нижних колонтитулов

%Пакеты шрифтов, кодировок. НЕ МЕНЯТЬ РАСПОЛОЖЕНИЕ.
\usepackage[utf8]{inputenc} % кодировка символов.
%\usepackage{mathtext} % позволяет использовать русские буквы в формулах. НЕСОВМЕСТИМО С tempora.
\usepackage[T1, T2A]{fontenc} % кодировка шрифта.
\usepackage[english, russian]{babel} % доступные языки.


%tikz:
\usepackage{tikz} % tikz
\usetikzlibrary{decorations.text} % позволяет делать текст вдоль кривой.
\usetikzlibrary{external} % позволяет кэшировать рисунки tikz.

%Отступы и поля:
%размеры страницы А4 11.7x8.3in
\textwidth=7.3in % ширина текста
\textheight=10in % высота текста
\oddsidemargin=-0.5in % левый отступ(базовый 1дюйм + значение)
\topmargin=-0.5in % отступ сверху до колонтитула(базовый 1дюйм + значение)


%Сокращения
%Скобочки
\newcommand{\inrad}[1]{\left( #1 \right)}
\newcommand{\inner}[1]{\left( #1 \right)}
\newcommand{\infig}[1]{\left{ #1 \right}}
\newcommand{\insqr}[1]{\left[ #1 \right]}
\newcommand{\ave}[1]{\left\langle #1 \right\rangle}


%% Красивые <= и >=
\renewcommand{\geq}{\geqslant}
\renewcommand{\leq}{\leqslant}

%% Более привычные греческие буквы
\renewcommand{\phi}{\varphi}
\renewcommand{\epsilon}{\varepsilon}
\newcommand{\eps}{\varepsilon}
\newcommand{\com}{\mathbb{C}}
\newcommand{\re}{\mathbb{R}}
\newcommand{\nat}{\mathbb{N}}

\newcommand{\stp}{$\filledmedtriangleleft$}
\newcommand{\enp}{$\filledmedsquare$}

\makeatletter
\newcommand{\sqbox}{%
    \collectbox{%
        \@tempdima=\dimexpr\width-\totalheight\relax
        \ifdim\@tempdima<\z@
            \fbox{\hbox{\hspace{-.5\@tempdima}\BOXCONTENT\hspace{-.5\@tempdima}}}%
        \else
            \ht\collectedbox=\dimexpr\ht\collectedbox+.5\@tempdima\relax
            \dp\collectedbox=\dimexpr\dp\collectedbox+.5\@tempdima\relax
            \fbox{\BOXCONTENT}%
        \fi
    }%
}
\makeatother
\newcommand{\mergelines}[2]{
\begin{tabular}{llp{.5\textwidth}}
#1 \\ #2
\end{tabular}
}
\newcommand\tab[1][0.51cm]{\hspace*{#1}}
\newcommand\difh[2]{\frac{\partial #1}{\partial #2}}
\newcommand{\messageforpeople}[1]{HSE Faculty of Physics \ \ HSE Faculty of Physics HSE Faculty of Physics \ \ HSE Faculty of Physics HSE Faculty of Physics \ \ HSE Faculty of Physics HSE Faculty of Physics \ \ HSE Faculty of Physics HSE Faculty of Physics \ \ HSE Faculty of Physics HSE Faculty of Physics \ \ HSE Faculty of Physics HSE Faculty of Physics \ \ HSE Faculty of Physics HSE Faculty of Physics \ \ HSE Faculty of Physics }


\numberwithin{equation}{section}

\begin{document}

%\begin{titlepage}
%    \pagestyle{empty}
%    \oddsidemargin=-1.18in
%    \topmargin=-1.5in
%    \begin{tikzpicture}
    \draw [white] (0, 0) ellipse (4.1in and 5.7in);
    \path [postaction={
    decoration={text along path,
    text={\messageforpeople{}},
    text align=center,
    raise=0mm,
    reverse path}, decorate}]
    [postaction={
    decoration={text along path,
    text={\messageforpeople{}},
    text align=center,
    raise=3mm,
    reverse path}, decorate}]
    [postaction={
    decoration={text along path,
    text={\messageforpeople{}},
    text align=center,
    raise=6mm,
    reverse path}, decorate}]
    [postaction={
    decoration={text along path,
    text={\messageforpeople{}},
    text align=center,
    raise=9mm,
    reverse path}, decorate}]
    [postaction={
    decoration={text along path,
    text={\messageforpeople{}},
    text align=center,
    raise=12mm,
    reverse path}, decorate}] 
    [postaction={
    decoration={text along path,
    text={\messageforpeople{}},
    text align=center,
    raise=15mm,
    reverse path}, decorate}]
    [postaction={
    decoration={text along path,
    text={\messageforpeople{}},
    text align=center,
    raise=18mm,
    reverse path}, decorate}]
    [postaction={
    decoration={text along path,
    text={\messageforpeople{}},
    text align=center,
    raise=21mm,
    reverse path}, decorate}] 
    [postaction={
    decoration={text along path,
    text={\messageforpeople{}},
    text align=center,
    raise=24mm,
    reverse path}, decorate}]
    [postaction={
    decoration={text along path,
    text={\messageforpeople{}},
    text align=center,
    raise=27mm,
    reverse path}, decorate}]
    [postaction={
    decoration={text along path,
    text={\messageforpeople{}},
    text align=center,
    raise=30mm,
    reverse path}, decorate}] 
    [postaction={
    decoration={text along path,
    text={\messageforpeople{}},
    text align=center,
    raise=33mm,
    reverse path}, decorate}]
    [postaction={
    decoration={text along path,
    text={\messageforpeople{}},
    text align=center,
    raise=36mm,
    reverse path}, decorate}]
    [postaction={
    decoration={text along path,
    text={\messageforpeople{}},
    text align=center,
    raise=39mm,
    reverse path}, decorate}] 
    [postaction={
    decoration={text along path,
    text={\messageforpeople{}},
    text align=center,
    raise=42mm,
    reverse path}, decorate}]
    [postaction={
    decoration={text along path,
    text={\messageforpeople{}},
    text align=center,
    raise=45mm,
    reverse path}, decorate}]
    [postaction={
    decoration={text along path,
    text={\messageforpeople{}},
    text align=center,
    raise=48mm,
    reverse path}, decorate}] 
    [postaction={
    decoration={text along path,
    text={\messageforpeople{}},
    text align=center,
    raise=51mm,
    reverse path}, decorate}]
    [postaction={
    decoration={text along path,
    text={\messageforpeople{}},
    text align=center,
    raise=54mm,
    reverse path}, decorate}]
    [postaction={
    decoration={text along path,
    text={\messageforpeople{}},
    text align=center,
    raise=57mm,
    reverse path}, decorate}] 
    [postaction={
    decoration={text along path,
    text={\messageforpeople{}},
    text align=center,
    raise=60mm,
    reverse path}, decorate}]
    [postaction={
    decoration={text along path,
    text={\messageforpeople{}},
    text align=center,
    raise=63mm,
    reverse path}, decorate}]
    [postaction={
    decoration={text along path,
    text={\messageforpeople{}},
    text align=center,
    raise=66mm,
    reverse path}, decorate}] 
    [postaction={
    decoration={text along path,
    text={\messageforpeople{}},
    text align=center,
    raise=69mm,
    reverse path}, decorate}]
    [postaction={
    decoration={text along path,
    text={\messageforpeople{}},
    text align=center,
    raise=72mm,
    reverse path}, decorate}]
    (0, 0) ellipse (2.75in and 4.96in);
    \node at (0, 3.2in) {\Large\textbf{Билеты по квантовой механике}};
    \node at (0, 0in) {\includegraphics[height=4in]{FOPF.png}};
    \node at (0, -3in) {Авторы:};
    \node at (0, -3.2in) {Карибджанов Матвей};
    \node at (0, -3.5in) {Januar 2023};
\end{tikzpicture}

%2.75in and 4.96in
%\end{titlepage}

\oddsidemargin=-0.5in
\topmargin=-0.5in
\pagestyle{plain}

\newpage
\tikzexternalize
\tableofcontents
\newpagestyle{main}{
\setfootrule{0.4pt}
\setfoot{}{\thepage}{Билет №\thesection}}
\pagestyle{main}

\section{Билет}
\subsection{соотношение неопределенности}
Говрит o том, что неовзмжно точно померить и координату и импульс.
 Так как если велечины измеримы то $\exists$ такая есть общий базис 
 из собствнных функций. Тогда если велечинам соответствуют 
 операторы $\hat A,\ \hat B$ то:
\begin{equation*} 
 \hat A \hat B \ket{n} = \hat B \hat A  \ket{n} \ \implies\ 
 \inner{\hat B \hat A - \hat A \hat B} \ket{n} = 
 \insqr{\hat A, \hat B}\ket{n} = 0,
\end{equation*}
а следовательно 
\begin{equation} \label{eq:1}
 \insqr{\hat A, \hat B} = \hat 0.
\end{equation} 
Согласн принципу неопределенности
\begin{equation} 
 \Delta_{\psi} A \Delta_{\psi} B \  
\end{equation} 
Например, для импульса и координаты ссотношение неопредленности 
выглядит следущим образом:
\begin{equation} 
 \Delta x \cdot \Delta p \geq \cfrac{\hslash}{2}.
\end{equation} 
Для доказательства этого соотношения представим, что
\begin{equation} 
 \insqr{\hat f, \hat g} = - i \hslash \hat C,
\end{equation}  
Теперь давайте в ведем:
\begin{equation} 
  \sigma \hat f = \sqrt{\ave{\inner{\hat f - \ave{\hat f}}^2 }} = \sqrt{\ave{\inner{ \ave{\hat f_0}^2 }}},
\end{equation} 
\begin{equation} 
 \ket{\phi } = \inner{\alpha \hat f + i \hat g }\ket{\psi}.
\end{equation} 
Так же не забываем требование эрмитовости операторов тогда:
\begin{equation} 
 \braket{\phi}{\phi} = \bra{\psi} \inner{\alpha^2 \hat f_0^2 + \hat g_0^2 + i\alpha \insqr{\hat f_0, \hat g_0}} \ket{\psi} = 
 \alpha^2 \ave{\hat f_0^2} +  \ave{ \hat g_0^2} + \alpha \hslash \ave{\hat C}
\end{equation} 

\subsection{Волновая функция}
Это удовлетворяет уравнению:
\begin{equation} 
 p (x, t) = \abs{\psi (x, t)}^2.
\end{equation} 
Для нее выполнятся условие нормировки 
\begin{equation} 
 \braket{\psi (x)}{\psi (x)} = \int_{-\infty}^{\infty} 
  \abs{\psi (x)}^2 dx = 1.
\end{equation} 
   
\subsection{Принцип суперпозиции}
Пусть состояне частици описыватся волновой функций $\psi_1$
 и втоже время $\psi_2$ приэтом каждой функции соответственно удовлетворят
 состояния $A_1, \ A_2$ то также эта частица будет описыватся волновой 
 волновой функцией $\psi = c_1 \psi_1 + c_2 \psi_2$ где $c_1, \ c_2 \in 
 \com$ c соответствующим состоянем $A$ это сотояние при 
 измерении может принять значения $A = A_1 \lor A = A_2$.

Пример частица может находитсяв точке $X_1$ этому полодение соответствует 
 собственная функция $\psi_1$, но также она модет быть в точке $X_2 - \psi_2$ 
 Тогда если мыбудем пытаться измерить полодение чатици то мыбудем получать
 значения $x_1 \lor x_2$, приэтом вероятность: 
 \begin{equation} 
  \mathbb{P}(X = x_1) = \braket{\psi}{\psi_1}.
 \end{equation}

 










\section{Билет}
\subsection{Операторы}
\subsubsection{Определение}
Оператор должен удовлетворять основному свойству
\begin{equation} 
 \hat A f(x) = \xi (x), 
\end{equation} 
То-есть оператор переводит Функцию в функцию, для 
любых двух операторов выполняется свойство асоциативности, 
но не для любых операторов выполнятся свойство комуативности. 
Как видно из \ref{eq:1}{уравнениея 1} из комутативности следует 
одновремнная измеримость.
\subsubsection{Основное свойство}
Оператор линеен
\begin{equation} 
  \hat A \inner{c_1 \ket{\psi_1} + c_2 \ket{\psi_2}} = 
  c_1 \hat A \ket{\psi_1} + c_2 \hat A \ket{\psi_2}
\end{equation} 
\subsubsection{Транспонирование}
\begin{equation} 
 \bra{\psi} \hat A \ket{\varphi} = \bra{\varphi} \hat A^T \ket{\psi} 
\end{equation}

\subsection{Связь измеримых велечин и операторов}
Во-первых, $\psi_n$ (обозначаем как $n$) собственная функция 
$\hat A $ если выполняется:
\begin{equation} 
 \hat A \ket{\psi_n} = A_n \ket{\psi_n}
\end{equation} 
При этом $A_n \in \com$ это величина, но что бы она была измерима 
надо чтобы  $A_n \in \re$, то условие выполняется в случае если: 
\begin{equation} 
    \hat A^\dag : = \hat A^{T *}.
\end{equation} 

\begin{equation} 
    \hat A^\dag \leqslant \hat A.
\end{equation} 













\end{document}